\chapter{Introduction générale}
\label{ch:intro}

\section{Contexte et motivation}

La plaidoirie est un acte oral préparé à l'écrit. Dans beaucoup de cabinets marocains, le circuit reste artisanal : notes manuscrites, dictée au collaborateur, ressaisie, mise en page Word peu avant l'audience. Ce circuit est lent, source d'erreurs de copie, et peu traçable.

Les logiciels de gestion d'affaires structurent le dossier ; les outils de dictée bureautique transcrivent hors SaaS ; les chatbots d'IA grand public n'offrent ni isolation multi-locataires, ni rôles avocat / assistant, ni export d'acte. LexHub vise cet écart : un \textbf{système d'information de cabinet} (SaaS) dont le métier maître est le \textbf{dossier + la plaidoirie}. L'IA (STT, OCR, LLM) est un \textbf{service partagé}, jamais le maître des faits\footnote{L'avocat demeure seul responsable du contenu plaidé ou déposé. LexHub assiste la rédaction ; il ne se substitue pas au conseil juridique.}. Le détail marché et la cartographie MADI figurent aux chapitres~\ref{ch:etat-art} et~\ref{ch:conception}.

Le stage PFA s'est déroulé du 1\textsuperscript{er} juillet au 31 août 2026 chez \textbf{WEBVOICE} (Tanger), sous tutelle de l'\textbf{ENSA Tétouan} (Université Abdelmalek Essaâdi). Encadrant entreprise : M.~Hatim Benhaddou. Tuteur pédagogique : Prof.~Mohammed Al Achhab. Le livrable est \textbf{LexHub}, en production à l'adresse \url{https://lexhub.space}.

\section{Présentation de l'organisme d'accueil --- WebVoice}
\label{sec:webvoice}

\textbf{WebVoice} (\url{https://webvoice.pro}) est une société de Tanger (29, rue Amr Ibn Ass, étage 3, n$^\circ$26). Gérant : M.~Khalid Merini Laarossi. Encadrant du stage : M.~Hatim Benhaddou (\texttt{hatim@webvoice.pro}).

L'entreprise se positionne sur la \textbf{voix sur IP (VoIP)}, le \textbf{développement web} et les \textbf{applications mobiles}. Sa démarche projet --- analyse, concept, mise en œuvre, suivi --- a structuré le PFA : cadrage du cœur plaidoirie, conception de l'espace cabinet, implémentation incrémentale, déploiement en production.

\subsection{Pourquoi ce sujet chez WebVoice}

WebVoice a identifié un besoin concret : un outil SaaS capable de transformer une dictée en document de plaidoirie exploitable, dans un cadre multi-cabinets. Le stage visait un produit \textbf{utilisable en production}, couvrant notamment l'authentification et les rôles, l'atelier de plaidoirie, l'équipe et le modèle d'abonnement (chapitres~\ref{ch:analyse} et~\ref{ch:economie}).

\section{Problématique}
\label{sec:problematique}

\paragraph{Constat.}
Préparer une plaidoirie suppose de relier la parole de l'avocat, les pièces du dossier et un document Word (.docx) prêt à relire ou à confier. Ces sources sont aujourd'hui dispersées.

\paragraph{Écart.}
Les logiciels de gestion n'offrent pas d'atelier vocal ; la dictée bureautique n'est pas un SaaS de cabinet ; les assistants d'IA grand public n'isolent pas les données et risquent d'\textbf{inventer} des faits.

\paragraph{Contraintes.}
Secret professionnel ; bilingue français / arabe ; l'assistant ne doit pas voir les brouillons ; l'IA ne doit pas inventer le fond ; le système doit être déployé.

\paragraph{Question.}
Comment concevoir, réaliser et mettre en production un SaaS multi-cabinets dans lequel l'avocat mène un atelier en \textbf{cinq étapes}, collabore avec un assistant limité aux \textbf{plaidoiries validées}, et dispose d'un modèle d'abonnement --- sans que l'IA n'invente le contenu juridique ?

\section{Objectifs du projet}

\begin{enumerate}[label=\textbf{O\arabic*.}]
  \item Atelier de plaidoirie déployé (dossier et pièces, dictée, correction, rédaction assistée, export Word).
  \item Espace cabinet : dossiers, documents, planning, activité, équipe.
  \item Rôles distincts : avocat, assistant, admin cabinet, admin plateforme.
  \item Validation comme seuil de visibilité pour l'assistant.
  \item Modèle économique (chapitre~\ref{ch:economie}) et exploitation en production.
\end{enumerate}

\section{Méthodologie de travail}

Approche agile, itérative, selon la démarche WebVoice (analyse, conception, mise en œuvre, suivi). Chaque itération a compris un cadrage avec les encadrants, la conception, le développement, des tests par rôle, puis un déploiement. Le rapport décrit exclusivement la version en ligne sur \url{https://lexhub.space}.

\section{Plan du rapport}

\begin{enumerate}[leftmargin=1.4em]
  \item Chapitre~\ref{ch:etat-art} --- marché, écart et positionnement LexHub.
  \item Chapitre~\ref{ch:analyse} --- acteurs, cas d'utilisation, exigences.
  \item Chapitre~\ref{ch:conception} --- architecture, urbanisation MADI, UML.
  \item Chapitre~\ref{ch:economie} --- modèle économique et rentabilisation.
  \item Chapitre~\ref{ch:realisation} --- réalisation et guide d'utilisation.
  \item Chapitre~\ref{ch:tests} --- tests, CI/CD, limites et perspectives.
  \item Chapitre~\ref{ch:conclusion} --- conclusion générale.
\end{enumerate}
